\ParallelTitle{A slower walk through a familiar street}{Redécouvrir une rue en marchant lentement}

\ParallelSection{notice}{Make room for noticing}{Prendre le temps de regarder}

\ParallelText{opening}{A familiar street can become little more than the space between two appointments. Try walking it once without an errand to finish. The aim is not to find an impressive landmark or collect a perfect photograph, but to notice details that a hurried journey usually edits out.}{Une rue familière finit parfois par n’être que le trajet entre deux rendez-vous. Essayez de la parcourir une fois sans course à terminer. Le but n’est ni de trouver un monument remarquable ni de réussir une photographie parfaite, mais de remarquer ce qu’un passage pressé laisse de côté.}

\ParallelText{change}{A painted door, a patch of shade or the sound of dishes from an open window can be enough. You do not need to explain every detail immediately. A useful observation may begin as a question, and a quiet walk can leave that question open until you have time to investigate it.}{Une porte peinte, un coin d’ombre ou le bruit de la vaisselle derrière une fenêtre ouverte peuvent suffire. Il n’est pas nécessaire de tout expliquer sur-le-champ. Une observation utile peut commencer par une question, que la promenade laisse ouverte jusqu’au moment où vous pourrez l’explorer.}

\ParallelSection{route}{Choose a modest route}{Choisir un itinéraire simple}

\ParallelText{route-plan}{Choose a starting point you can reach easily and an ending point that does not require you to rush back. A short loop is often more useful than an ambitious itinerary. Check the actual conditions where you walk; the small map below is an invented illustration, not directions for a real place.}{Choisissez un départ facile d’accès et une arrivée qui ne vous oblige pas à rentrer en hâte. Une petite boucle est souvent plus utile qu’un programme ambitieux. Vérifiez les conditions sur place : le schéma ci-dessous est imaginaire et ne donne pas d’indications pour un lieu réel.}

\ParallelFigure{route-image}{images/figure-route-image.png}{An invented route with four stopping points. The numbers identify places to pause; they do not indicate distance, time or accessibility.}{Un parcours imaginaire avec quatre haltes. Les numéros désignent des arrêts, sans indiquer leur distance, leur durée ni leur accessibilité.}

\ParallelText{steps.item-1}{\begin{itemize}\item \phantomsection\label{steps}\hypertarget{steps}{}Pick one detail at each stop rather than trying to describe everything around you.\end{itemize}}{\begin{itemize}\item \ifdefstring{\ParallelMode}{right}{\phantomsection\label{steps}\hypertarget{steps}{}}{}Choisissez un détail à chaque halte, au lieu de tout décrire autour de vous.\end{itemize}}

\ParallelText{steps.item-2}{\begin{itemize}\item Leave a little unplanned time. If something deserves another look, you can stay.\end{itemize}}{\begin{itemize}\item Gardez du temps libre. Vous pourrez vous attarder si quelque chose mérite un second regard.\end{itemize}}

\ParallelText{steps.item-3}{\begin{itemize}\item Keep other people’s privacy in mind. A written observation does not require a photograph of a stranger.\end{itemize}}{\begin{itemize}\item Respectez la vie privée des autres. Une observation écrite n’exige pas de photographier un inconnu.\end{itemize}}

\clearpage

\ParallelSection{notebook}{Keep a small notebook}{Tenir un petit carnet}

\ParallelText{notebook-prose}{A notebook can hold fragments rather than polished sentences. Record what you actually noticed before adding an interpretation. “The bench was empty” is an observation; “nobody uses this square” is a much broader claim. Keeping those two kinds of statement separate makes later writing more honest and more interesting.}{Le carnet peut accueillir des fragments, sans phrases parfaitement rédigées. Notez d’abord ce que vous avez réellement observé, avant de l’interpréter. « Le banc était vide » est une observation ; « personne ne fréquente cette place » est une affirmation bien plus large. Les distinguer rendra votre récit plus juste et plus intéressant.}

\ParallelText{prompt}{\begin{quote}One sound, one colour, one question: a simple prompt for your own notes.\end{quote}}{\begin{quote}Un son, une couleur, une question : une consigne simple pour vos notes.\end{quote}}

\begin{ParallelKeep}

\ParallelText{notes-table.row-0}{\phantomsection\label{notes-table}\hypertarget{notes-table}{}\begin{tabularx}{\linewidth}{@{}>{\RaggedRight\arraybackslash}X>{\RaggedRight\arraybackslash}X@{}}\toprule \textbf{Observation} & \textbf{Question}\\ \midrule \end{tabularx}}{\ifdefstring{\ParallelMode}{right}{\phantomsection\label{notes-table}\hypertarget{notes-table}{}}{}\begin{tabularx}{\linewidth}{@{}>{\RaggedRight\arraybackslash}X>{\RaggedRight\arraybackslash}X@{}}\toprule \textbf{Observation} & \textbf{Question}\\ \midrule \end{tabularx}}

\ParallelText{notes-table.row-1}{\begin{tabularx}{\linewidth}{@{}>{\RaggedRight\arraybackslash}X>{\RaggedRight\arraybackslash}X@{}}An empty bench & Is it used at another time?\\ \end{tabularx}}{\begin{tabularx}{\linewidth}{@{}>{\RaggedRight\arraybackslash}X>{\RaggedRight\arraybackslash}X@{}}Un banc vide & Est-il occupé à une autre heure ?\\ \end{tabularx}}

\ParallelText{notes-table.row-2}{\begin{tabularx}{\linewidth}{@{}>{\RaggedRight\arraybackslash}X>{\RaggedRight\arraybackslash}X@{}}A newly painted door & What colour was it before?\\ \end{tabularx}}{\begin{tabularx}{\linewidth}{@{}>{\RaggedRight\arraybackslash}X>{\RaggedRight\arraybackslash}X@{}}Une porte repeinte & Quelle était sa couleur avant ?\\ \end{tabularx}}

\ParallelText{notes-table.row-3}{\begin{tabularx}{\linewidth}{@{}>{\RaggedRight\arraybackslash}X>{\RaggedRight\arraybackslash}X@{}}Music from a window & Does the sound change farther away?\\ \bottomrule \end{tabularx}}{\begin{tabularx}{\linewidth}{@{}>{\RaggedRight\arraybackslash}X>{\RaggedRight\arraybackslash}X@{}}De la musique à une fenêtre & Le son change-t-il en s’éloignant ?\\ \bottomrule \end{tabularx}}

\ParallelText{notes-table.caption}{These invented notes show the difference between an observation and a question. They are examples, not findings about a real neighbourhood.}{Ces notes imaginaires distinguent observation et question. Ce sont des exemples, pas des constats sur un quartier réel.}\end{ParallelKeep}

\ParallelSection{return}{Read the walk again}{Relire la promenade}

\ParallelText{return-prose}{After returning, choose one note and expand it into a short paragraph. Keep the observed detail, explain why it mattered to you and mark anything you still do not know. You can check a factual question later, but you need not turn every walk into a research project. Sometimes the value lies in having paid attention.}{Au retour, choisissez une note et développez-la en un court paragraphe. Conservez le détail observé, expliquez pourquoi il vous a touché et signalez ce que vous ignorez encore. Vous pourrez vérifier un fait plus tard, sans transformer chaque promenade en enquête. Parfois, sa valeur tient simplement à l’attention que vous lui avez portée.}

\ParallelText{route-reference}{\ParallelReference{route}{Return to the route section}}{\ParallelReference{route}{Revenir à la partie sur l’itinéraire}}
